% Generated from locale/tr/content/methods/induction/induction-on-N.tex; do not hand-edit.


\olfileid[tr]{mth}{ind}{inN}

\olsection{$\Nat$ Üzerinde Tümevarım}

En basit biçimiyle tümevarım, bütün doğal sayılar için geçerli
sonuçları kanıtlamakta kullanılan bir tekniktir. $0$'dan başlayıp
tekrar tekrar~$1$ ekleyerek sonunda her doğal sayıya ulaşmamız olgusunu
kullanır. Dolayısıyla bir şeyin her sayı için doğru olduğunu kanıtlamak
üzere (1)~$0$ için doğru olduğunu belirleyebilir ve (2)~bir~$n$ sayısı
için doğru olduğunda sonraki~$n+1$ sayısı için de doğru olduğunu
gösterebiliriz. ``~$n$ sayısı $P$ özelliğine sahiptir'' sözünü $P(n)$
ile (ve ``~$k$ sayısı $P$ özelliğine sahiptir'' sözünü $P(k)$ ile vb.)
kısaltırsak, bütün $n \in \Nat$ için $P(n)$'nin tümevarımla kanıtı
şunlardan oluşur:
\begin{enumerate}
\item $P(0)$'ın bir kanıtı ve
\item herhangi bir~$k$ için, $P(k)$ ise $P(k+1)$ olduğunun bir kanıtı.
\end{enumerate}
Bunu bütünüyle açık kılmak için hem (1)'in hem de~(2)'nin elimizde
olduğunu varsayın. O zaman (1), $P(0)$'ın doğru olduğunu söyler. (2) de
elimizdeyse özellikle $P(0)$ ise $P(0+1)$, yani $P(1)$ olduğunu biliriz.
Bu, ``herhangi bir~$k$ için, $P(k)$ ise $P(k+1)$'' genel ifadesinde~$k$
yerine~$0$ konarak çıkar. Dolayısıyla modus ponens ile~$P(1)$'i elde
ederiz. Yeniden (2)'den, bu kez~$k$ yerine $1$ alarak, $P(1)$ ise
$P(2)$ olduğunu elde ederiz. $P(1)$'i az önce belirlediğimizden modus
ponens ile~$P(2)$'yi elde ederiz. Bu böyle sürer. Herhangi bir~$n$
sayısı için bunu $n$~kez yaptıktan sonra sonunda~$P(n)$'ye ulaşırız.
Öyleyse (1) ve~(2) birlikte herhangi $n \in \Nat$ için~$P(n)$'yi
belirler.

Bir örneğe bakalım. $n$ zarla kaç farklı toplam atabileceğimizi bulmak
istediğimizi düşünün. Saçma görünebilir ama $0$ zarla başlayalım. Hiç
zarınız yoksa ``atabileceğiniz'' yalnızca bir olası toplam vardır: hiç
benek yoktur ve toplam~$0$'dır. Dolayısıyla farklı olası atışların
sayısı~$1$'dir. Yalnızca bir zarınız varsa, yani $n=1$ ise, $1$'den~$6$'ya
altı olası değer vardır. İki zarla $2$'den $12$'ye kadar her toplamı
atabiliriz; bu $11$ olasılıktır. Üç zarla $3$'ten $18$'e kadar her
sayıyı, yani $16$ farklı olasılığı atabiliriz. $1$, $6$, $11$, $16$:
  bir örüntüye benziyor; yanıt $5n+1$ olabilir mi? Elbette en fazla
  $5n+1$ olası değer vardır; çünkü $n$ zarla atabileceğiniz en küçük değer olan
$n$ (bütün zarlar $1$) ile atabileceğiniz en büyük değer olan $6n$
(bütün zarlar $6$) arasında yalnızca $5n+1$ sayı vardır.

\begin{thm}
  $n$ zarla, $n$ ile $6n$ arasındaki $5n+1$ olası değerin tamamı
  atılabilir.
\end{thm}

\begin{proof}
$P(n)$ şu iddia olsun: ``$n$~zar kullanarak $n$ ile~$6n$ arasındaki
herhangi bir sayıyı atmak mümkündür.'' Tümevarımı kullanmak için şunları
kanıtlarız:
\begin{enumerate}
\item \emph{Tümevarım tabanı} $P(1)$, yani yalnızca bir zarla $1$ ile
  $6$ arasındaki herhangi bir sayıyı atabilirsiniz.
\item \emph{Tümevarım adımı}: bütün $k$'ler için, $P(k)$ ise~$P(k+1)$.
\end{enumerate}

(1), $6$ yüzlü bir zar incelenerek kanıtlanır. Altı yüzünün tamamı
vardır ve $1$ ile~$6$ arasındaki her sayı yüzlerden birinde yer alır.
Dolayısıyla tek bir zarla $1$ ile~$6$ arasındaki herhangi bir sayıyı
atmak mümkündür.

(2)'yi kanıtlamak için koşullunun önbileşenini, yani $P(k)$'yi
varsayarız. Bu varsayıma \emph{tümevarım hipotezi} denir. Onu $P(k+1)$'i
kanıtlamakta kullanırız. Zor olan, $k+1$ zarın olası atış değerlerini
  $k$~zarın olası atış değerleriyle ve ek $k+1$'inci zarın atışlarıyla
ilişkilendirerek düşünmenin bir yolunu bulmaktır---tümevarım hipotezini
kullanmak istiyorsak yapmamız gereken budur.

Tümevarım hipotezi, $k$~zar kullanarak $k$ ile~$6k$ arasındaki herhangi
bir sayıyı elde edebileceğimizi söyler. $(k+1)$'inci zarımızla~$1$
atarsak bu, toplama $1$ ekler. Böylece $k$~zar atıp ardından
$(k+1)$'inci zarla~$1$ atarak $k+1$ ile $6k+1$ arasındaki herhangi bir
değeri atabiliriz. Geriye ne kalır? $6k+2$'den $6k+6$'ya kadar olan
değerler. Bunları $k$ tane $6$ atıp ardından $(k+1)$'inci zarımızla $2$
ile $6$ arasında bir sayı atarak elde edebiliriz. Bunlar birlikte,
$k+1$ zarla $k+1$ ile $6(k+1)$ arasındaki sayıların herhangi birini
atabileceğimiz anlamına gelir; yani tümevarım hipotezi olan~$P(k)$
varsayımını kullanarak~$P(k+1)$'i kanıtladık.
\end{proof}

Çoğu kez, kendisi ``tümevarımsal'' olarak tanımlanan, yani $(n+1)$'inci
nesnesi $n$'inci nesne cinsinden tanımlanan bir nesneler dizisi
(sayılar, kümeler vb.) hakkında bir şey kanıtlamak istediğimizde
tümevarımı kullanırız. Örneğin $n$'ye kadar olan doğal sayıların
toplamı~$s_n$'yi şöyle tanımlayabiliriz:
\begin{align*}
  s_0 & = 0\\
  s_{n+1} & = s_n + (n+1)
\end{align*}
Bu tanım şunları verir:
\begin{align*}
  s_0 & = 0,\\
  s_1 & = s_0 + 1 && = 1,\\
  s_2 & = s_1 + 2 && = 1 + 2 = 3\\
  s_3 & = s_2 + 3 && = 1 + 2 + 3 = 6, \text{ vb.}
\end{align*}
Artık $s_n = n(n+1)/2$ olduğunu tümevarımla kanıtlayabiliriz.

\begin{prop}
  $s_n = n(n+1)/2$.
\end{prop}

\begin{proof}
  (1) $s_0 = 0\cdot(0 + 1)/2$ olduğunu ve (2) $s_k = k(k+1)/2$ ise
  $s_{k+1} = (k+1)(k+2)/2$ olduğunu kanıtlamalıyız. (1) açıktır.
  (2)'yi kanıtlamak için tümevarım hipotezini varsayarız: $s_k =
  k(k+1)/2$. Bunu kullanarak $s_{k+1} = (k+1)(k+2)/2$ olduğunu
  göstermeliyiz.

  $s_{k+1}$ nedir? Tanım gereği $s_{k+1} = s_k + (k+1)$. Tümevarım
  hipotezi gereği $s_k = k(k+1)/2$. Bunu bir önceki eşitlikte yerine
  koyabiliriz; ardından yalnızca biraz kesir aritmetiğine ihtiyacımız
  kalır:
  \begin{align*}
    s_{k+1} & = \frac{k(k+1)}{2} + (k+1) = {}\\
    & = \frac{k(k+1)}{2} + \frac{2(k+1)}{2} = {}\\
    & = \frac{k(k+1) + 2(k+1)}{2} = {}\\
    & = \frac{(k+2)(k+1)}{2}.
  \end{align*}
\end{proof}

Buradaki önemli ders şudur: tümevarımsal olarak tanımlanmış bir $a_n$
dizisi hakkında bir şey kanıtlıyorsanız, tümevarım başvurulacak en açık
yoldur. Açık yol olmasa bile (zar atışlarının olasılıklarında olduğu
gibi),~$k+1$ durumunu~$k$ durumuyla bir biçimde ilişkilendirebilirseniz
tümevarımı kullanabilirsiniz.

